\section{Limitations and Ethical Considerations}

This is a single-benchmark development study, not formal-test confirmation. MOCHEG development outcomes were already inspected in earlier work, so the expanded validation check is reused development evidence rather than statistically fresh confirmation. The official test remains sealed. No current second benchmark supports an unchanged replication; MMCV would require a table modality and a newly defined split, while AVerImaTeC would require a new evidence-identity and open-web corpus protocol.

The image system uses filename text but no OCR or generated caption, because those sources were outside the frozen trial budget. Text reranking is expensive (about 19 minutes for the full reranker pass on an RTX 4090) and degrades early precision. The membership models are low-capacity logistic controls rather than fine-tuned multimodal cross-encoders. The expected-F1 decoder assumes conditional independence, although evidence from one article or duplicate group is correlated. Source metadata are insufficient for a source-cluster bootstrap, so claims are the resampling unit.

The 2,189 claims without a released non-empty complete-gold mapping are excluded transparently. This may limit population coverage, but treating them as true empty evidence sets would reward empty predictions and conflate release mapping with task semantics. Controlled robustness is explicitly post-outcome and exploratory. In particular, the positive distractor-promotion result is an artificial stress test and cannot support a deployment claim.

Automated fact-checking may misrepresent people, sources, or events when retrieval is incomplete. The system should support human evidence review, not autonomously adjudicate truth. Public data and model artifacts are used under their research licenses; immutable identifiers and hashes are retained, while credentials and private information are not distributed. Negative results and failed gates are preserved to reduce selective reporting.
